\documentclass[10pt,a4paper]{article}
\usepackage[T1]{fontenc}
\usepackage{lmodern}
\usepackage[margin=1.8cm]{geometry}
\usepackage{amsmath,booktabs}
\usepackage{microtype}
\usepackage[colorlinks=true,linkcolor=black,urlcolor=blue!50!black]{hyperref}
\hypersetup{pdftitle={Backtest Overfitting Statistics: Research Summary},pdfauthor={Mayank Khandelwal}}
\usepackage{enumitem}
\setlist{nosep,leftmargin=1.2em}
\pagestyle{plain}

\begin{document}
\begin{center}
{\large\textbf{Backtest Overfitting Statistics under Serial Dependence}}\\[2pt]
{A Computational Replication and Empirical Stress Test: research summary}\\[2pt]
{\small Mayank Khandelwal \quad Working paper, 12 November 2026}
\end{center}

\paragraph{Question.} PSR, the 2014 DSR and CSCV/PBO are often read as hurdles (``DSR $>0.95$'', ``PBO $\le0.05$''). How do they behave as decision rules when returns are serially dependent, heteroskedastic, heavy tailed or cross-correlated, and when a historical search winner is tested against a real benchmark?

\paragraph{What was done.} Independent, test-driven reimplementation of PSR, MinTRL, DSR, MinBTL and CSCV/PBO with page-level provenance; selected numerical examples from the four source papers reproduced. Zero-edge Monte Carlo (2000 datasets per cell) under IID, AR(1), GARCH(1,1), Student-$t$ and equicorrelated returns. Circular block bootstrap of daily Kenneth French 49-industry returns, 1971--2026, centred to zero edge, including a sub-period split at 1999. A search over 93 moving-average rules on the Fama--French market factor (1972--2026, 5 bps per position change), tested against buy-and-hold with a paired HAC Sharpe-difference test, Bonferroni and Hansen's SPA, plus exploratory lag, cost, sub-period and walk-forward checks. Every saved output was regenerated from the current code and matched the archived results.

\paragraph{Main findings.}
\begin{enumerate}
\item \textbf{PSR and dependence (confirms known theory).} At zero Sharpe, uncorrected PSR(0)$>0.95$ rejects 11.85\% [10.47, 13.35] under AR(1) $\phi=0.3$; the large-sample prediction $1-\Phi(1.645/\sqrt F)$ with $F=(1+\phi)/(1-\phi)$ gives 11.37\%. GARCH, Student-$t$ and correlated designs stay at 5\% (as expected at $\mathrm{SR}^\ast=0$). A 7-lag Newey--West variance reduces the AR(1) rate to 7.10\%, not 5\%.
\item \textbf{PSR on real data (new evidence, conditional).} With each period centred, resampled 1971--1998 industry returns give 8.35\% [7.17, 9.65] and 1999--2026 give 3.85\% [3.05, 4.79]: the error runs in both directions and tracks lag-1 autocorrelation ($+0.157$ vs $-0.026$ on average). A long-run-variance approximation lands within 0.1 points in both periods. The full-window rate (about 6\%) is a mixture.
\item \textbf{The 2014 DSR is not a 5\% test.} On the best of 100 IID zero-skill strategies it rejects 0.05\% [0.00, 0.28] (PSR: 99.15\%). Acceptance rises with cross-strategy correlation (2.60\% at $\rho=0.7$; 5.05\% with the Eq.~9 implied $N=30.7$), because $V$ omits the common component.
\item \textbf{PBO under exchangeable nulls.} Mean PBO stays near 1/2 (or 25/49) in every zero-edge design; such nulls do not test PBO's power, which was not measured.
\item \textbf{Case study.} Selected rule 15/225: Sharpe 0.588 vs market 0.458, DSR 0.998, PBO 0.208. Paired HAC difference $+0.130$ (s.e.\ 0.099, one-sided $p=0.094$); 0 of 93 rules pass Bonferroni; SPA $p=0.504$. By a Normal approximation holding the s.e.\ fixed, 80\% power needs a difference of 0.246 (0.406 with Bonferroni). Walk-forward re-selection (exploratory) earns 0.466 vs the market's 0.536 over 1982--2026 (corrected for switch costs).
\end{enumerate}

\begin{center}\small
\begin{tabular}{@{}lll@{}}
\toprule
Decision & Statistic & Does not establish \\
\midrule
One pre-specified Sharpe above a hurdle & PSR & Robustness to serial dependence \\
Search winner better than best of $N$ noise trials & DSR (2014) & A 5\% error rate; skill vs a benchmark \\
Rankings persist out of sample & PBO & A probability of future loss \\
Move capital from benchmark to strategy & none of these & needs a search-aware benchmark test with power \\
\bottomrule
\end{tabular}
\end{center}

\paragraph{Contribution and novelty limits.} Replication with traceable provenance; quantified operating characteristics of the 2014 DSR and of PSR on real daily data; an explicit contrast between search-adjusted zero-skill inference and benchmark-relative inference, with power. Not claimed: a new statistic, first identification of the serial-dependence problem (Lo 2002; Ledoit and Wolf 2008), or any profitable strategy. The 2026 DSR framework of L\'opez de Prado and Porcu, examined only through its abstract, may subsume the DSR findings.

\paragraph{Status of evidence.} Designs and predictions were recorded before runs; from the case study on they were committed to version control before running. The evidence rule, the sub-period prediction and all robustness checks were fixed after earlier results were known and are labelled exploratory or post hoc. One earlier sub-period result was withdrawn after an audit found the periods were not centred.

\paragraph{Main limitations.} One data vendor and market; $\mathrm{SR}^\ast=0$ only for PSR; fixed bootstrap block lengths; idealized same-close execution on a research portfolio; HAC assumptions unverified for 54 years of data; external data access/redistribution terms; both input archives now have timestamped, hash-verified recovery URLs.

\end{document}
